\section{Objetivos}

\begin{itemize}
      \item \textbf{Mapear e quantificar as taxas de depreciação dos veículos elétricos:} Mensurar a velocidade e amplitude da perda de valor dos modelos eletrificados em comparação direta a veículos a combustão pareados por porte, categoria e faixa de preço de lançamento \cite{schloter2022empirical}.
      
      \item \textbf{Avaliar o efeito de transbordamento (\textit{spillover}) de marcas entrantes:} Analisar correlações temporais e possíveis quebras estruturais de preços em segmentos tradicionais após a introdução e guerra tarifária promovida por montadoras chinesas no mercado nacional \cite{li2026global, li2025rise}.

      \item \textbf{Implementar e validar o \textit{reshaping} de séries temporais em imagens 2D:} Transformar as trajetórias unidimensionais de preços de cada veículo em representações matriciais utilizando algoritmos como \textit{Gramian Angular Field} (GAF) ou \textit{Markov Transition Field} (MTF) \cite{wu2023imaging, jueco2024performance}.

      \item \textbf{Identificar e caracterizar arquétipos de desvalorização veicular:} Aplicar algoritmos de agrupamento (\textit{clustering}) sobre representações vetoriais clássicas e sobre os espaços latentes gerados a partir das imagens, identificando perfis comportamentais típicos (alta retenção, desvalorização acelerada, valorização tardia).

      \item \textbf{Comparar empiricamente a eficácia das representações:} Avaliar quantitativamente, por meio de métricas de qualidade de agrupamento (como coeficiente de Silhueta e índice Davies-Bouldin), se a transformação das séries para imagem supera as abordagens vetoriais unidimensionais tradicionais, adotando a distância Euclidiana e o \textit{Dynamic Time Warping} (DTW) como \textit{baselines} de comparação.
\end{itemize}
